\documentclass[10pt,a4paper]{amsart}
\usepackage[margin=19mm]{geometry}
\usepackage{amsmath,amssymb,mathtools}
\usepackage[scr=rsfs,cal=euler]{mathalfa}
\usepackage{xfrac}
\usepackage[unicode,hidelinks,bookmarks=false]{hyperref}
\newcommand{\ie}{i.e.\ }
\newcommand{\cf}{cf.\ }
\newcommand{\resp}{respectively\ }
\newcommand{\Subsection}{\subsection}
\providecommand{\nobreakdash}{\nobreak\hbox{-}}
\setlength{\emergencystretch}{2em}
\multlinegap0pt
\usepackage{bm}
\usepackage{bbm}
\usepackage{dsfont}
\usepackage{xspace}
\usepackage{cancel}
\usepackage{mathtools}
\usepackage[shortlabels]{enumitem}
\usepackage{amsthm}
\makeatletter
\@ifundefined{theorem}{\newtheorem{theorem}{Theorem}}{}
\@ifundefined{definition}{\newtheorem{definition}[theorem]{Definition}}{}
\@ifundefined{lemma}{\newtheorem{lemma}[theorem]{Lemma}}{}
\makeatother
\newcommand{\N}{\mathbb{N}}
\newcommand{\Pcal}{\mathcal{P}}
\newcommand{\T}{\mathbb{T}}
\newcommand{\dd}{\,\mathrm{d}}
\title[Statistical Error Bounds for Generative Solvers of Chaotic P]{Statistical Error Bounds for Generative Solvers of Chaotic PDEs: Wasserstein Stability, Generalization, and Turbulence}
\author{Victor Armegioiu}
\date{Source: arXiv:2602.18794v1 (CC BY 4.0)}
\begin{document}
\maketitle

\noindent\emph{Source and licence.} Statistical Error Bounds for Generative Solvers of Chaotic PDEs: Wasserstein Stability, Generalization, and Turbulence. Version arXiv:2602.18794v1 (CC BY 4.0), distributed under the Creative Commons Attribution 4.0 International licence, as recorded in the arXiv licence metadata for this exact version and shown on the abstract page https://arxiv.org/abs/2602.18794v1. Full rights notice: licenses/arxiv-2602.18794-CC-BY-4.0.txt.\par\smallskip

\noindent\textit{Editorial scope.} Counted unit: the Lemma whose source label is \texttt{lem:time-reg-closed} in the article (source0.tex lines 389--400, source label \texttt{lem:time-reg-closed}), delivered together with the article's own complete original proof of it (source0.tex lines 402--529, one \texttt{proof} environment of 128 lines). No mathematics is written, repaired or re-derived: statement and proof are the article's own text line for line. Required context retained uncounted: def:time-regular (lines 371--381). Every definition, display and label of those lines is retained unchanged. Omitted from this package and not redistributed: 1--370, 382--388, 401--401, 530--3900. Nothing inside a retained excerpt is omitted or altered: every retained body is delivered exactly as the article's lines are written, so no omission receipt of any kind accompanies this package. Citations: the retained excerpts carry no citation command at all, so no bibliography entry of the article is transcribed and no reference is redistributed. Reference adaptation: none was needed and none was made; every reference inside the delivered unit resolves inside it, so no unresolved reference or citation is produced. Printed slips kept literal: no printed word, symbol, hypothesis or number of the counted statement or of its proof was altered. Layout and class: the article's own class is \texttt{article} with packages including enumitem, geometry, fontenc, inputenc, lmodern, babel, microtype, xcolor, mathtools, amsmath, amssymb, amsthm, bbm, graphicx; the wrapper is the lane's pinned \texttt{amsart} editorial wrapper at 10pt on A4, which restates the theorem-like environments the retained text needs and adds the article's own macro definitions that the retained text needs (\texttt{\textbackslash N}, \texttt{\textbackslash Pcal}, \texttt{\textbackslash T}, \texttt{\textbackslash dd}), with extra wrapper packages (bm, bbm, dsfont, xspace, cancel, mathtools, enumitem). The delivered numbering is the unit's own. Assets: no retained span contains a figure environment, a graphics-inclusion command, a PDF-page inclusion, a table or a TikZ picture; no image file is copied into this package, the package declares no asset dependency and no image file is redistributed with it. Licence: version arXiv:2602.18794v1 of this article is distributed under the Creative Commons Attribution 4.0 International licence, as recorded in the arXiv licence metadata for this exact version and shown on the abstract page https://arxiv.org/abs/2602.18794v1. A full rights notice is delivered at licenses/arxiv-2602.18794-CC-BY-4.0.txt. Characters: every retained excerpt is pure ASCII, so the delivered engine, XeLaTeX with its default Latin Modern text font, reports no missing character.
\begin{definition}[Time-regularity]\label{def:time-regular}
A curve $\mu_\cdot\in L^1_t(\Pcal)$ is \emph{time-regular} if there exist $L\in\N$, $C>0$, and a measurable assignment
$(s,t)\mapsto\pi_{s,t}\in\Pcal(L^2_x\times L^2_x)$ such that for a.e.\ $s,t\in[0,T)$:
\begin{enumerate}[label=(\roman*),leftmargin=2.2em]
\item $\pi_{s,t}\in\Pi(\mu_s,\mu_t)$;
\item \begin{equation}\label{eq:time-reg}
\int_{L^2_x\times L^2_x}\|u-v\|_{H^{-L}}\,\dd\pi_{s,t}(u,v)\le C|t-s|.
\end{equation}
\end{enumerate}
A family $\{\mu^\Delta_\cdot\}_{\Delta>0}$ is \emph{uniformly time-regular} if the same $(C,L)$ works for all $\Delta$.
\end{definition}

\begin{lemma}[LM time-regularity is closed under $d_T$ limits]
\label{lem:time-reg-closed}
Let $\mu^m_\cdot\in L^1([0,T);\Pcal(L^2_x))$ be time-regular with the same constants $(C,L)$ in the sense of
Definition~\ref{def:time-regular}. Assume
\[
d_T(\mu^m,\mu)\to 0
\quad\text{and}\quad
\sup_m\int_0^T\int_{L^2_x}\|u\|_2\,\dd\mu^m_t(u)\,\dd t<\infty.
\]
Assume $D=\T^d$ so that the embedding $L^2_x\hookrightarrow H^{-L}(D)$ is compact.
Then $\mu_\cdot$ is time-regular with the same $(C,L)$.
\end{lemma}

\begin{proof}
Let $\lambda:=\frac{1}{T^2}\,\dd s\,\dd t$ on $[0,T]^2$.
For each $m$, let $(s,t)\mapsto\pi^m_{s,t}\in\Pi(\mu^m_s,\mu^m_t)$ be a measurable assignment such that
\[
\int_{L^2_x\times L^2_x}\|u-v\|_{H^{-L}}\,\dd\pi^m_{s,t}(u,v)\le C|t-s|
\quad\text{for a.e.\ }(s,t)\in[0,T]^2.
\]

\smallskip
\noindent\textbf{Step 1: build averaged couplings and extract a limit in $H^{-L}$.}
Define $\Theta^m\in\Pcal([0,T]^2\times L^2_x\times L^2_x)$ by
\[
\Theta^m(\dd s\,\dd t\,\dd u\,\dd v)
:=\lambda(\dd s\,\dd t)\,\pi^m_{s,t}(\dd u,\dd v).
\]
View $L^2_x\times L^2_x$ as embedded in $H^{-L}\times H^{-L}$.
Let $B_R:=\{u\in L^2_x:\|u\|_2\le R\}$. Since $B_R$ is relatively compact in $H^{-L}$ and
\[
\Theta^m(\{(s,t,u,v):\|u\|_2>R\})
=\frac1T\int_0^T \mu^m_s(\{\|u\|_2>R\})\,\dd s
\le \frac{1}{RT}\int_0^T\int\|u\|_2\,\dd\mu^m_s\,\dd s,
\]
(and similarly for $v$), the family $\{\Theta^m\}$ is tight on $[0,T]^2\times H^{-L}\times H^{-L}$.
Hence, after extracting a subsequence (not relabeled),
\[
\Theta^m\Rightarrow\Theta
\quad\text{weakly in }\Pcal\big([0,T]^2\times H^{-L}\times H^{-L}\big).
\]
The $(s,t)$-marginal of each $\Theta^m$ is $\lambda$, hence the $(s,t)$-marginal of $\Theta$ is also $\lambda$.
Disintegrate $\Theta$ w.r.t.\ $\lambda$:
there exists a $\lambda$-a.e.\ defined measurable family $(s,t)\mapsto\pi_{s,t}\in\Pcal(H^{-L}\times H^{-L})$ such that
\[
\Theta(\dd s\,\dd t\,\dd u\,\dd v)=\lambda(\dd s\,\dd t)\,\pi_{s,t}(\dd u,\dd v).
\]

\smallskip
\noindent\textbf{Step 2: identify the marginals using $H^{-L}$-continuous tests.}
Let $\iota:L^2_x\to H^{-L}$ denote the continuous embedding and set
\[
\mu^{m,H}_t:=\iota_\#\mu^m_t\in\Pcal(H^{-L}),
\qquad
\mu^{H}_t:=\iota_\#\mu_t\in\Pcal(H^{-L}).
\]
Since $d_T(\mu^m,\mu)=\int_0^T W_1(\mu^m_t,\mu_t)\,\dd t\to0$, after extracting a further subsequence we may assume
\[
W_1(\mu^m_t,\mu_t)\to0 \quad\text{for a.e.\ }t\in[0,T].
\]
For such $t$, $W_1(\mu^m_t,\mu_t)\to0$ implies $\mu^m_t\Rightarrow\mu_t$ narrowly in $L^2_x$, and therefore (by continuity of $\iota$)
\[
\mu^{m,H}_t=\iota_\#\mu^m_t \Rightarrow \iota_\#\mu_t=\mu^H_t
\quad\text{narrowly in }H^{-L},\ \text{for a.e.\ }t.
\]

Fix $\psi\in C([0,T]^2)$ and $\varphi\in C_b(H^{-L})$.
Define $F(s,t,u,v):=\psi(s,t)\varphi(u)$, which is bounded continuous on $[0,T]^2\times H^{-L}\times H^{-L}$.
Then $\Theta^m\Rightarrow\Theta$ gives
\begin{equation}\label{eq:step2-theta-limit}
\int F\,\dd\Theta^m \longrightarrow \int F\,\dd\Theta.
\end{equation}
Compute the left-hand side using $\Theta^m=\lambda\otimes\pi^m_{s,t}$ and the fact that the first marginal of $\pi^m_{s,t}$ is $\mu^m_s$:
\[
\int F\,\dd\Theta^m
=
\int_{[0,T]^2}\psi(s,t)\Big(\int_{L^2_x}\varphi(\iota(u))\,\dd\mu^m_s(u)\Big)\dd\lambda(s,t)
=
\int_{[0,T]^2}\psi(s,t)\Big(\int_{H^{-L}}\varphi(u)\,\dd\mu^{m,H}_s(u)\Big)\dd\lambda(s,t).
\]
For a.e.\ $s$, $\mu^{m,H}_s\Rightarrow\mu^H_s$, hence $\int\varphi\,\dd\mu^{m,H}_s\to\int\varphi\,\dd\mu^H_s$.
Since $|\int\varphi\,\dd\mu^{m,H}_s|\le\|\varphi\|_\infty$, dominated convergence yields
\[
\int F\,\dd\Theta^m
\longrightarrow
\int_{[0,T]^2}\psi(s,t)\Big(\int_{H^{-L}}\varphi(u)\,\dd\mu^{H}_s(u)\Big)\dd\lambda(s,t).
\]
On the other hand, disintegration of $\Theta$ gives
\[
\int F\,\dd\Theta
=
\int_{[0,T]^2}\psi(s,t)\Big(\int_{H^{-L}\times H^{-L}}\varphi(u)\,\dd\pi_{s,t}(u,v)\Big)\dd\lambda(s,t).
\]
Comparing with \eqref{eq:step2-theta-limit} and using arbitrariness of $\psi$, we conclude that for $\lambda$-a.e.\ $(s,t)$,
\[
\int_{H^{-L}\times H^{-L}}\varphi(u)\,\dd\pi_{s,t}(u,v)=\int_{H^{-L}}\varphi(u)\,\dd\mu^H_s(u)
\qquad\forall \varphi\in C_b(H^{-L}),
\]
i.e.\ the first marginal of $\pi_{s,t}$ is $\mu^H_s$. The same argument with $F(s,t,u,v)=\psi(s,t)\varphi(v)$ shows that
the second marginal is $\mu^H_t$.

Since $\mu^H_s,\mu^H_t$ are supported on $\iota(L^2_x)$, $\pi_{s,t}$ is supported on $\iota(L^2_x)\times\iota(L^2_x)$.
Because $\iota$ is continuous and injective between Polish spaces, Lusin--Souslin implies $\iota(L^2_x)$ is Borel in $H^{-L}$
and $\iota^{-1}:\iota(L^2_x)\to L^2_x$ is Borel. Define
\[
\widetilde\pi_{s,t}:=(\iota^{-1}\times \iota^{-1})_\#\pi_{s,t}\in\Pcal(L^2_x\times L^2_x).
\]
Then $\widetilde\pi_{s,t}\in\Pi(\mu_s,\mu_t)$ for $\lambda$-a.e.\ $(s,t)$.

\smallskip
\noindent\textbf{Step 3: pass the increment bound to the limit.}
Fix $\psi\ge0$ in $C([0,T]^2)$ and set $G(s,t,u,v):=\psi(s,t)\|u-v\|_{H^{-L}}$ on $[0,T]^2\times H^{-L}\times H^{-L}$.
Then $G\ge0$ is continuous, hence by Portmanteau,
\[
\int G\,\dd\Theta
\le \liminf_{m\to\infty}\int G\,\dd\Theta^m.
\]
But for each $m$,
\[
\int G\,\dd\Theta^m
=
\int_{[0,T]^2}\psi(s,t)\Big(\int\|u-v\|_{H^{-L}}\,\dd\pi^m_{s,t}\Big)\dd\lambda(s,t)
\le C\int_{[0,T]^2}\psi(s,t)|t-s|\,\dd\lambda(s,t).
\]
Therefore,
\[
\int G\,\dd\Theta
\le C\int_{[0,T]^2}\psi(s,t)|t-s|\,\dd\lambda(s,t).
\]
Disintegrating $\Theta$ yields that for $\lambda$-a.e.\ $(s,t)$,
\[
\int_{H^{-L}\times H^{-L}}\|u-v\|_{H^{-L}}\,\dd\pi_{s,t}(u,v)\le C|t-s|.
\]
Since $\iota$ is the canonical injection of $L^2_x$ into $H^{-L}(D)$ (identifying $L^2$ functions with distributions),
we have $\|\iota(u)-\iota(v)\|_{H^{-L}}=\|u-v\|_{H^{-L}}$ for all $u,v\in L^2_x$. Therefore, pushing forward by $\iota^{-1}\times\iota^{-1}$ yields, for $\lambda$-a.e.\ $(s,t)$,
\[
\int_{L^2_x\times L^2_x}\|u-v\|_{H^{-L}}\,\dd\widetilde\pi_{s,t}(u,v)\le C|t-s|.
\]

Thus $\mu_\cdot$ is time-regular with constants $(C,L)$.
\end{proof}

\end{document}
